% Part: normal-modal-logic
% Chapter: filtrations
% Section: more-filtrations

\documentclass[../../../include/open-logic-section]{subfiles}

\begin{document}

\olfileid[tr]{nml}{fil}{acc}

\olsection{Süzmeler ve Erişilebilirlik Özellikleri}

Belirtildiği gibi, evrensel, seri ve yansımalı modellerin süzmeleri de her zaman
sırasıyla evrensel, seri ya da yansımalıdır. Buna karşılık simetrik veya geçişli
bir modelin her süzmesi sırasıyla simetrik veya geçişli olmak zorunda değildir.
Yine de bazı durumlarda süzmeleri bu özellikleri koruyacak biçimde tanımlamak
mümkündür. Bunun için en kaba süzmenin tanımındaki gibi ilerler, fakat $R^*$
tanımına ek koşullar koyarız. $\Gamma$ alt-!!{formula}s altında kapalı olsun.
$\mModel{M}=\tuple{W,R,V}$ modelindeki $u$ ve $v$ dünyaları arasında
\olref{tab:Cn-filtrations}'de verilen $C_i(u,v)$ bağıntılarını düşünün.
$R^*[u][v]$'yi bu koşulların birleşimlerine dayanarak tanımlayabiliriz.
Örneğin $R^*[u][v]$ ancak ve ancak $C_1(u,v)$ geçerliyse koşulunu koyarsak
tam olarak en kaba süzmeyi elde ederiz. $R^*[u][v]$ ancak ve ancak hem
$C_1(u,v)$ hem $C_2(u,v)$ geçerliyse dersek başka bir süzme elde ederiz.
Bu süzme en kaba süzmeden ``daha incedir''; çünkü $C_1(u,v)$ ile
$C_2(u,v)$'yi birlikte sağlayan dünya çiftleri, yalnızca $C_1(u,v)$'yi
sağlayan çiftlerin bir alt kümesidir.

\begin{table}[ht]
  \centering
  \begin{tabular}{|ll|}
    \hline
    \multirow{2}{*}{$C_1(u,v)$:}
    \iftag{prvBox}{&
      $\Box!A \in \Gamma$ ve $\mSat{M}{\Box!A}[u]$ ise
      $\mSat{M}{!A}[v]$;\iftag{prvDiamond}{ ve}{}\\}{}
    \iftag{prvDiamond}{&
      $\Diamond!A \in \Gamma$ ve $\mSat{M}{!A}[v]$ ise
      $\mSat{M}{\Diamond!A}[u]$; \\}{}
    \hline
    \multirow{2}{*}{$C_2(u,v)$:}
    \iftag{prvBox}{&
      $\Box!A \in \Gamma$ ve $\mSat{M}{\Box!A}[v]$ ise
      $\mSat{M}{!A}[u]$;\iftag{prvDiamond}{ ve}{}\\}{}
    \iftag{prvDiamond}{&
      $\Diamond!A \in \Gamma$ ve $\mSat{M}{!A}[u]$ ise
      $\mSat{M}{\Diamond!A}[v]$; \\}{}
    \hline
    \multirow{2}{*}{$C_3(u,v)$:}
    \iftag{prvBox}{&
      $\Box!A \in \Gamma$ ve $\mSat{M}{\Box!A}[u]$ ise
      $\mSat{M}{\Box!A}[v]$;\iftag{prvDiamond}{ ve}{}\\}{}
    \iftag{prvDiamond}{&
      $\Diamond!A \in \Gamma$ ve $\mSat{M}{\Diamond!A}[v]$ ise
      $\mSat{M}{\Diamond!A}[u]$; \\}{}
    \hline
    \multirow{2}{*}{$C_4(u,v)$:}
    \iftag{prvBox}{&
      $\Box!A \in \Gamma$ ve $\mSat{M}{\Box!A}[v]$ ise
      $\mSat{M}{\Box!A}[u]$;\iftag{prvDiamond}{ ve}{}\\}{}
    \iftag{prvDiamond}{&
      $\Diamond!A \in \Gamma$ ve $\mSat{M}{\Diamond!A}[u]$ ise
      $\mSat{M}{\Diamond!A}[v]$; \\}{}
    \hline
    \end{tabular}
  \caption{Süzmeleri tanımlamak için olası dünyalara konan koşullar.}
  \ollabel{tab:Cn-filtrations}
\end{table}

\begin{thm}\ollabel{thm:more-filtrations}
  $\mModel{M}=\tuple{W,R,V}$ bir model, $\Gamma$ alt-!!{formula}s altında
  kapalı olsun. $W^*$ ile $V^*$, \olref[fil]{defn:filtration}'deki gibi
  tanımlansın. O zaman:
  \begin{enumerate}
  \item $R^*[u][v]$ ancak ve ancak $C_1(u,v)\land C_2(u,v)$ olsun.
    Bu durumda $R^*$ simetriktir ve $\mModel{M}$ simetrikse
    $\mModel{M^*}=\tuple{W^*,R^*,V^*}$ bir süzmedir.
  \item $R^*[u][v]$ ancak ve ancak $C_1(u,v)\land C_3(u,v)$ olsun.
    Bu durumda $R^*$ geçişlidir ve $\mModel{M}$ geçişliyse
    $\mModel{M^*}=\tuple{W^*,R^*,V^*}$ bir süzmedir.
  \item $R^*[u][v]$ ancak ve ancak $C_1(u,v)\land C_2(u,v)
    \land C_3(u,v)\land C_4(u,v)$ olsun. Bu durumda $R^*$ simetrik ve
    geçişlidir; $\mModel{M}$ simetrik ve geçişliyse
    $\mModel{M^*}=\tuple{W^*,R^*,V^*}$ bir süzmedir.
  \item $R^*[u][v]$ ancak ve ancak $C_1(u,v)\land C_3(u,v)
    \land C_4(u,v)$ olacak biçimde $R^*$ tanımlansın. Bu durumda $R^*$
    geçişli ve Öklidyendir; $\mModel{M}$ geçişli ve Öklidyense
    $\mModel{M^*}=\tuple{W^*,R^*,V^*}$ bir süzmedir.
  \end{enumerate}
\end{thm}

\begin{proof}
  \begin{enumerate}
    \item $C_1(u,v)\Leftrightarrow C_2(v,u)$ ve
      $C_2(u,v)\Leftrightarrow C_1(v,u)$ olduğundan $R^*$'nin simetrik
      olduğu doğrudan görülür. Dolayısıyla yalnızca $\mModel{M}$ simetrikse
      $\mModel{M^*}$'nin $\Gamma$ üzerinden bir süzme olduğunu göstermemiz
      gerekir. $C_1(u,v)$ koşulu,
      \iftag{prvBox}{%
        \iftag{prvDiamond}{%
          \olref[fil]{defn:filtration-R2} ile
          \olref[fil]{defn:filtration-R3} koşullarının
          \olref[fil]{defn:filtration} içinde}{%
          \olref[fil]{defn:filtration}\olref[fil]{defn:filtration-R2}
          koşulunun}}{%
        \olref[fil]{defn:filtration}\olref[fil]{defn:filtration-R3}
        koşulunun}
      sağlanmasını güvence altına alır. Bu nedenle yalnızca
      \olref[fil]{defn:filtration}\olref[fil]{defn:filtration-R1}'i,
      yani $Ruv$'nin $R^*[u][v]$'yi gerektirdiğini doğrulamalıyız.

      Öyleyse $Ruv$ olsun. $R^*[u][v]$'yi göstermek için $C_1(u,v)$ ile
      $C_2(u,v)$'yi kurmamız gerekir. $C_1$ için:
      \iftag{prvBox}{$\Box!A\in\Gamma$ ve
        $\mSat{M}{\Box!A}[u]$ ise $Ruv$ olduğundan
        $\mSat{M}{!A}[v]$ de geçerlidir.\iftag{prvDiamond}{ Benzer
        biçimde, }{}}{}
      \iftag{prvDiamond}{$\Diamond!A\in\Gamma$ ve
        $\mSat{M}{!A}[v]$ ise $Ruv$ olduğundan
        $\mSat{M}{\Diamond!A}[u]$.}{}
      $C_2$ için:
      \iftag{prvBox}{$\Box!A\in\Gamma$ ve
        $\mSat{M}{\Box!A}[v]$ ise simetriklik ve $Ruv$, $Rvu$'yu
        gerektirir; dolayısıyla $\mSat{M}{!A}[u]$.
        \iftag{prvDiamond}{ Benzer biçimde, }{}}{}
      \iftag{prvDiamond}{$\Diamond!A\in\Gamma$ ve
        $\mSat{M}{!A}[u]$ ise simetriklikten $Rvu$, dolayısıyla
        $\mSat{M}{\Diamond!A}[v]$.}{}
    \item Alıştırma.
    \item Alıştırma.
    \item Alıştırma.
  \end{enumerate}
\end{proof}

\begin{prob}
  \olref[nml][fil][acc]{thm:more-filtrations} teoreminin ispatını
  tamamlayın.
\end{prob}

\end{document}
